% main.tex — CLEI 2026 paper (Intelligent Systems track)
\documentclass[conference]{IEEEtran}

\usepackage[T1]{fontenc}
\usepackage{cite}
\usepackage{graphicx}
\usepackage{amsmath,amssymb}
\usepackage{booktabs}
\usepackage{algpseudocode}
\usepackage{float}
\usepackage{url}

% Define algorithm float (compatible with IEEEtran)
\floatstyle{ruled}
\newfloat{algorithm}{tbp}{loa}
\floatname{algorithm}{Algorithm}

\begin{document}

\title{Early Population Dynamics Provide Stronger Switching Signals than Static Landscape Features for Warmup-Based Operator Switching}

\author{\IEEEauthorblockN{Pedro Sanches Zambrano,
Eduardo Faria de Souza,
and Celso G. Camilo-Junior}
\IEEEauthorblockA{Institute of Informatics, Federal University of Goi\'as\\
Goi\^ania, GO, Brazil\\
pedrosanches@discente.ufg.br}}

\maketitle

%% ===================================================================
\begin{abstract}
Adaptive operator control often relies on reward signals accumulated
during the run, yet for some operators a single keep-or-drop decision
after a short warmup may suffice. We investigate whether static
fitness-landscape features or early population dynamics better inform
such a decision for an IVF-inspired intensification operator inside
SPEA2. On 51 continuous multi-objective benchmark instances, a
14-feature landscape feature set provides moderate signal for ranking
expected IGD gain (Random Forest leave-one-out $r_s = 0.374$), but
binary classification does not generalize reliably across benchmark families.
In contrast, mean archive turnover during the first 20\% of generations
yields a stronger switching signal ($q_{\mathrm{BH}} = 0.013$,
$\hat{A}_{12} = 0.849$); a single-threshold rule achieves balanced
accuracy $0.754$ under leave-one-family-out cross-validation. A
warmup-based controller using this signal incurs no hypervolume
losses relative to always-on IVF-SPEA2 on the 10 IGD-labeled
\textsc{Not helpful} instances, while recording 34 wins against
baseline SPEA2 overall, at the cost of forfeiting gains on
9 \textsc{Helpful} cases concentrated in WFG.
Our results suggest that, for this operator--host pair, static
features serve as priors on gain magnitude whereas early dynamics are
more actionable for binary switching.
\end{abstract}

\begin{IEEEkeywords}
Fitness landscape analysis, Adaptive operator control,
Multi-objective optimization, Evolutionary algorithms, Population dynamics
\end{IEEEkeywords}

%% ===================================================================
\section{Introduction}
\label{sec:intro}

Whether a search operator helps a multi-objective evolutionary algorithm
depends on both the problem instance and the phase of the run. An
intensification step that accelerates convergence on one landscape may waste
evaluations on another.

Fitness landscape analysis (FLA) describes optimization problems through
numerical features of sampled search behavior and objective structure. In
single-objective continuous optimization, such descriptors have supported
algorithm selection with measurable gains over fixed
solvers~\cite{mersmann2011exploratory,kerschke2019automated,bischl2016aslib}.
Multi-objective extensions exist~\cite{liefooghe2018landscape}, but their
empirical use for operator-level decisions is limited.

Most FLA studies target full algorithm selection. Here we study a narrower
decision: after a warmup phase, should one intensification operator remain
active inside a fixed host algorithm? The question is adjacent to adaptive
operator selection~\cite{fialho2008adaptive,maturana2009adaptive} and
component-wise algorithm design~\cite{bezerra2016automatic}, but we focus on
one operator, one host algorithm, and one irreversible keep/drop decision
rather than repeated on-line adaptation.

We investigate this question for an IVF-inspired operator inside
SPEA2~\cite{zitzler2001spea2} on 51 instances from the ZDT, DTLZ, WFG, and
MaF families, examining static landscape features, early population dynamics,
and a one-shot warmup switch. Throughout the paper, \textsc{Helpful} and
\textsc{Not helpful} denote operational IGD-based analysis labels derived from
this specific benchmark cohort, not universal claims about the problem
instances.

Our contributions are:
\begin{enumerate}
	\item Static landscape descriptors provide moderate prior signal about
	      operator gain magnitude, but not robust cross-family binary
	      decisions: four retained features survive BH correction, Random
	      Forest regression reaches $r_s = 0.374$ and $R^{2} = 0.160$ in
	      LOOCV, and removing structural metadata weakens the signal.
	\item Early turnover is the clearest switching signal in this study,
	      whereas static binary prediction does not generalize across
	      benchmark families
	      ($q_{\mathrm{BH}} = 0.013$, $\hat{A}_{12} = 0.849$; threshold-rule
	      balanced accuracy $0.691$ in LOOCV and $0.754$ in leave-one-family-out
	      refitting, with uncertainty discussed in
	      Section~\ref{sec:early_turnover}).
	\item We evaluate a one-shot warmup controller based on early turnover,
	      showing that it avoids hypervolume losses relative to always-on
	      IVF-SPEA2 on all IGD-labeled \textsc{Not helpful} instances, while
	      forfeiting gains on 9~\textsc{Helpful} cases concentrated in
	      WFG (Section~\ref{sec:ctrl_results}).
\end{enumerate}

\subsection{Relationship to companion work}
The IVF operator and its empirical validation on this benchmark suite are
presented in a companion preprint~\cite{zambrano2026ivfspea2}, which addresses
the design, ablation, and tuning of the memetic hybrid. The companion preprint
evaluates whether the operator improves SPEA2; it does not study adaptive
control or landscape features. The present paper asks a different
question---whether static landscape features or early population dynamics
can inform a one-shot keep/drop decision for the operator---and contributes
all landscape analysis, dynamics experiments, threshold calibration, and
controller evaluation reported here. These reused IGD outcomes serve only to
define fixed external labels; no landscape feature extraction, trajectory
logging, threshold calibration, controller execution, or HV-based controller
evaluation is reused from the companion preprint. The two papers share a
benchmark suite and a host
algorithm, but address distinct research questions with distinct evidence. This paper should not be read as another performance evaluation of IVF/SPEA2, but as an operator-control study using IVF/SPEA2 as a fixed case study.

%% ===================================================================
\section{Background}
\label{sec:background}

\subsection{SPEA2}
\label{sec:spea2}

SPEA2~\cite{zitzler2001spea2} is an elitist multi-objective evolutionary
algorithm that maintains a population and an external archive of fixed size
$\bar{N}$. Fitness assignment combines a strength-based raw fitness with a
$k$-nearest-neighbor density estimate; environmental selection retains all
individuals with fitness below~1, applying distance-based truncation when the
non-dominated set exceeds $\bar{N}$. The full formulation appears
in~\cite{zitzler2001spea2}.

\subsection{The IVF-Inspired Crossover Operator}
\label{sec:ivf}

The intensification operator studied here derives from a family of
\emph{in vitro fertilization} (IVF) metaheuristics originally proposed for
single-objective search and later extended to multi-objective
settings~\cite{camilo-junior2011,sampaio2017ivf,sampaio2019ivf,Sampaio2024}.
The operator introduces two mechanisms:
\textbf{(H1)~Dissimilar father selection} selects a fraction $C$ of the
population as ``mothers'' and pairs each mother with a father chosen for high
objective-space distance (top-3 candidates, binary tournament on fitness).
\textbf{(H2)~Collective acceptance criterion} applies IVF cycles
(SBX,~$\eta_c = 20$)~\cite{deb1995simulated} up to
$\mathcal{C}_{\max}$ per generation while the \emph{average population
	fitness improves}, rather than requiring each offspring to outperform its
parent individually.
Algorithm~\ref{alg:ivf} gives the complete procedure.
The operator is skipped once its cumulative IVF evaluation count exceeds a
budget gate proportional to the current global evaluation count, controlled by
threshold $R$. The IVF budget gate $R$ controls cumulative IVF evaluations and
is independent of the controller warmup horizon used later for the keep/drop
decision (Section~\ref{sec:dynamics}). The broader algorithm design and empirical validation
appear in the companion preprint~\cite{zambrano2026ivfspea2}; unless stated otherwise, parameters are
$C = 0.12$, $R = 0.225$, mutation fraction
$p_{\mathrm{mut}} = 0.3$, variable mutation $V = 0.1$,
and $\mathcal{C}_{\max} = 2$.

We use the term IVF-inspired because the operator preserves the
collection-and-intensification logic of IVF but adapts the father-selection
and acceptance mechanisms to SPEA2's fitness assignment.

\begin{algorithm}[t]
\caption{IVF operator (one generation)}\label{alg:ivf}
\begin{algorithmic}[1]\footnotesize
\State \textbf{Input:} population $P$, fitness values $\phi$, params.\ $C, R, \mathcal{C}_{\max}$
\If{cum.\ IVF evals.\ $> R \times$ current global evals.}
  \State \Return $P$ unchanged \Comment{gate}
\EndIf
\State Sort $P$ by $\phi$ (ascending); mothers $\mathcal{M} \gets$ top $\lceil C \cdot N \rceil$
\State Father pool $\mathcal{P}_f \gets$ top 50\% of $P$ by $\phi$
\For{cycle $= 1, \dots, \mathcal{C}_{\max}$}
  \State $\bar{\phi}_{\mathrm{before}} \gets \mathrm{mean}(\phi)$
  \State $O \gets \emptyset$
  \For{each $m \in \mathcal{M}$} \Comment{H1: dissimilar father}
    \State Find 3 farthest fathers in $\mathcal{P}_f$ from $m$ (objective-space Euclidean distance)
    \State Pick father $p$ by binary tournament on $\phi$ among those 3
    \State $O \gets O \cup \{\mathrm{SBX}(m, p;\; \eta_c = 20)\}$
  \EndFor
  \State $P \gets P \cup O$; apply SPEA2 env.\ selection $\to$ size $N$
  \State Recompute fitness values $\phi$ for the new population
  \If{$\mathrm{mean}(\phi) \geq \bar{\phi}_{\mathrm{before}}$} \Comment{H2: collective crit.}
    \State \textbf{break} \Comment{no improvement}
  \EndIf
\EndFor
\State \Return $P$
\end{algorithmic}
\end{algorithm}

\subsection{Fitness Landscape Analysis for Multi-Objective Optimization}
\label{sec:fla_background}

Fitness landscape analysis characterizes optimization problems through
computable numerical features~\cite{mersmann2011exploratory}. For
single-objective problems, ELA features have supported effective algorithm
selection~\cite{kerschke2019automated}. The multi-objective setting is harder
because the landscape is vector-valued and must capture dominance structure,
Pareto front geometry, and objective correlation.

Liefooghe et al.~\cite{liefooghe2018landscape} proposed landscape features
for multi-objective optimization based on random walks and sample-based
probes, capturing distributional properties of objective values, dominance
relationships, and fitness-distance correlations. Whether such features carry
useful signal for operator-level decisions---rather than full algorithm
selection---has received little attention.


%% ===================================================================
\section{Can Static Features Predict Operator Benefit?}
\label{sec:fla}

\subsection{Experimental Setup}
\label{sec:fla_setup}

We evaluate IVF-SPEA2 and standard SPEA2 on 51 synthetic benchmark instances
drawn from the ZDT~\cite{zitzler2000comparison},
DTLZ~\cite{deb2005scalable}, WFG~\cite{huband2006review}, and
MaF~\cite{cheng2017benchmark} suites, with $M \in \{2, 3\}$ objectives.
Each synthetic instance contributes 60 independent IVF-SPEA2 runs and
60 independent SPEA2 runs, with a budget of $10^5$ function evaluations and
population size $N = 100$. All experiments are conducted on the
PlatEMO platform~\cite{tian2017platemo}.

One engineering case (RWMOP9) exists in the broader trajectory batch but is
excluded from the label-based analyses here.

IGD is the primary endpoint for label construction and static modeling. HV
remains the secondary endpoint in the synthetic host comparison and becomes
the primary endpoint in the controller section to avoid reusing the IGD label
metric. HV is computed on the final archive using a reference point of
$\mathbf{1}_M$ (vector of ones) after objective-wise normalization as
$(x - f_{\min})/[1.1(f_{\max} - f_{\min})]$, where $f_{\min}$ and
$f_{\max}$ are derived from the benchmark reference Pareto front
of each instance.
The same normalization and reference point are used for all algorithms
per instance.

The inferential core of this paper is the host comparison IVF-SPEA2 versus
SPEA2, since the response variable is the marginal gain of the IVF module
over its host. For broader algorithmic context, exploratory pairwise Wilcoxon comparisons
on the same 51-instance cohort show that IVF-SPEA2 wins 47/51 IGD cases
against NSGA-II~\cite{deb2002fast} and 40/51 against MOEA/D~\cite{zhang2007moead}, but only 26/51 against AR-MOEA~\cite{tian2018indicator}
and 26/51 against NSGA-III~\cite{deb2014evolutionary}, placing it in a competitive but not uniformly
dominant region. IVF-SPEA2 also incurs roughly $5{\times}$ runtime overhead
versus SPEA2 (median 60\,s vs.\ 11\,s per run at $10^5$ evaluations).
These contextual results motivate the operator-level question but do not
replace the confirmatory IVF-SPEA2 versus SPEA2 contrast.

Because outcome samples are independent (not seed-paired), we use the
Wilcoxon rank-sum (Mann--Whitney $U$) test and the Vargha--Delaney
$\hat{A}_{12}$ effect size~\cite{arcuri2011practical,vargha2000critique} on
the Inverted Generational Distance (IGD) of the final archive relative to
the true Pareto front.
For the binary analysis, an instance is labeled \textsc{Helpful} if IVF-SPEA2 is
better at $p < 0.05$ with $\hat{A}_{12} > 0.56$ (at least a small effect).
$\hat{A}_{12}$ is oriented so that values above~0.5 favor IVF-SPEA2, i.e.,
indicate lower IGD than SPEA2. All
remaining instances are labeled \textsc{Not helpful}. This yields
41~\textsc{Helpful} and 10~\textsc{Not helpful} in the binary view (the 10
comprise 6~neutral and 4~hurts cases). These labels are analysis categories
for this benchmark cohort, not family-wise corrected claims of universal
per-instance truth.

\subsection{Landscape Feature Extraction}
\label{sec:fla_features}

We extract 22 raw landscape features spanning five descriptor families:
objective statistics, dominance structure, front geometry, local-walk
behavior, and structural metadata. Representative examples include
objective correlation, skewness, dominance-depth variability,
front diameter, front-gap ratio, information content, neutrality ratio,
fitness--distance correlation, number of objectives, and number of
variables.

Features are computed from 2{,}000 Latin Hypercube sample points and
1{,}000-step probe walks, replicated across 5 independent seeds to reduce
sampling noise; seed-level estimates are arithmetically averaged before
modeling. All feature extraction operates on raw (unscaled) objective
values to preserve geometry-sensitive properties such as front diameter;
objective-wise $z$-score normalization is applied only during the PCA
projection in Fig.~\ref{fig:pca}. For modeling and visualization, we retain 14 features after
removing degenerate variables (single-value fraction $>0.90$) and
high-redundancy pairs (absolute Spearman correlation $|r| > 0.80$,
greedy drop by mean absolute correlation).
Because raw objective values preserve benchmark-specific scale, geometry
features such as front diameter may partially encode scale differences
across suites; normalized feature variants should be tested in future work.

\subsection{Predictive Modeling}
\label{sec:fla_prediction}

We first assess whether the retained feature sets predict operator
benefit in two settings:

\subsubsection{Regression}
The target variable is the rank-transformed median IGD difference
$\Delta\mathrm{IGD}$ between SPEA2 and IVF-SPEA2, where positive values
indicate IVF advantage.
A Random Forest~\cite{breiman2001random} with 500 trees is trained using leave-one-out
cross-validation (LOOCV). With the 14-feature retained set, the model
achieves Spearman $r_s = 0.374$ ($p = 0.0068$), $R^{2} = 0.160$, and
mean absolute error $= 11.00$ on the rank scale. Under
leave-one-family-out validation, the overall Spearman correlation remains
$0.380$, indicating that static descriptors do capture part of the
between-instance variation in operator gain.

\subsubsection{Classification}
For the binary label (\textsc{Helpful}/\textsc{Not helpful}), the Random
Forest achieves $\mathrm{BA} = 0.527$ ($\mathrm{MCC} = 0.059$) under LOOCV.
Alternative classifiers score higher under LOOCV (logistic regression
$\mathrm{BA} = 0.728$; RBF-SVC $\mathrm{BA} = 0.690$), but all collapse
to $\mathrm{BA} = 0.439$--$0.488$ under leave-one-family-out validation.
Static features help rank expected gain, but they do not support cross-family
binary switching. This negative classification result is central in this
study: static descriptors alone do not justify a reliable cross-family
keep/drop rule.

\subsection{Correlation Analysis}
\label{sec:fla_correlation}

We compute Spearman rank correlations between each retained feature and
$\Delta\mathrm{IGD}$, applying the Benjamini--Hochberg
procedure~\cite{benjamini1995controlling} at $q = 0.05$.

Four features survive correction (Table~\ref{tab:fla_top5}). The strongest
associations are front diameter and number of objectives (both positive),
followed by dominance-depth variability (negative) and number of variables
(positive).

\begin{table}[t]
	\centering
	\caption{Top five retained landscape features by $|\rho|$ (Spearman
		correlation with $\Delta\mathrm{IGD}$). Four features survive
		Benjamini--Hochberg correction at $q = 0.05$.}
	\label{tab:fla_top5}
	\begin{tabular}{llccc}
		\toprule
		\textbf{Group}                  & \textbf{Feature}                & $\boldsymbol{\rho}$ &
		$\boldsymbol{p_{\mathrm{raw}}}$ & $\boldsymbol{p_{\mathrm{adj}}}$                                         \\
		\midrule
		Front geometry                  & front\_diameter                 & $0.47$              & 0.0005 & 0.0037 \\
		Structural                      & n\_obj                          & $0.47$              & 0.0005 & 0.0037 \\
		Dominance                       & dom\_depth\_std                 & $-0.36$             & 0.0106 & 0.0380 \\
		Structural                      & n\_var                          & $0.35$              & 0.0109 & 0.0380 \\
		Local walk                      & fdc\_mean                       & $-0.21$             & 0.1368 & 0.3277 \\
		\bottomrule
	\end{tabular}
\end{table}

\subsection{Discussion}
\label{sec:fla_discussion}

Static descriptors carry signal about \emph{how much} IVF is likely to help,
but not enough for cross-family binary keep/drop decisions. Regression
captures part of the between-instance variation in $\Delta\mathrm{IGD}$ and
four features survive BH correction, yet the two classes overlap heavily in
feature space (Fig.~\ref{fig:pca}).

Part of that signal is carried by structural metadata (number of objectives,
number of variables). Removing them drops regression from $r_s = 0.374$ to
$0.314$ in LOOCV and from $0.380$ to $0.192$ in leave-one-family-out
validation; classification is weak in both cases.

The static FLA result should be read narrowly. The retained descriptors
contain useful prior information about gain magnitude, but the stronger LOOCV
classification signal disappears in family-held-out tests.
Other multi-objective landscape feature frameworks, such as PLON or Deep-ELA,
remain untested and could change this picture.

Ranking expected gain, however, is not sufficient for an in-run keep/drop
decision. The next question is whether within-run signals are sharper than
these static priors.
Fig.~\ref{fig:classifier} contrasts the cross-family classification
performance of static models against the early-turnover rule, confirming
that the stronger LOOCV signal does not generalize.

\begin{figure}[t]
	\centering
	\includegraphics[width=\columnwidth]{figures/fig1_pca_fla.pdf}
	\caption{PCA projection of the 14 retained landscape features
		for 51 benchmark instances, colored by operator benefit label.
		\textsc{Helpful} (circles) and \textsc{Not helpful} (triangles)
		overlap substantially. Despite the moderate continuous correlation
		between individual features and expected IGD gain, no linear
		combination cleanly separates the two classes.}
	\label{fig:pca}
\end{figure}

\begin{figure}[t]
	\centering
	\includegraphics[width=\columnwidth]{figures/fig7_classifier_comparison.pdf}
	\caption{Static versus dynamic signals for binary operator switching.
		Cross-family (LFO-CV) performance of static-landscape classifiers and
		the early-turnover rule. Static methods collapse below random
		(BA~=~0.5, dashed), while early turnover provides actionable signal
		(BA~=~0.754, MCC~=~0.413).}
	\label{fig:classifier}
\end{figure}


%% ===================================================================
\section{Population Dynamics and Early Signal}
\label{sec:dynamics}

Since static models do not generalize for cross-family binary selection, we
turn from ex ante descriptors to warmup summaries from paired trajectories.
The goal is to identify a signal observable during a run that
separates \textsc{Helpful} from \textsc{Not helpful} cases.

\subsection{Experimental Design}
\label{sec:dynamics_setup}

For the dynamic analysis, we execute 30 paired-seed runs on the same 51
labeled synthetic cases. IVF-SPEA2 and SPEA2 share the same random seed per
run, reducing variance due to stochastic initialization. Each run logs
per-generation snapshots including IGD$(t)$, HV$(t)$, spread$(t)$, spacing$(t)$, archive
turnover$(t)$ (fraction of archive members at generation $t$ that are replaced
relative to generation $t-1$), non-dominated fraction$(t)$,
and IVF cycle count$(t)$. This yields $51 \times 30 \times 2 = 3{,}060$
synthetic trajectory runs. The operator-benefit labels used throughout this
section and the next are those derived from the independent-run cohort
(60~runs per algorithm, Section~\ref{sec:fla_setup}); they are applied
unchanged to this paired-seed cohort. One additional engineering case was
logged in the broader batch but is excluded from the label-based tests.

\subsection{Trajectory Patterns and Self-Regulation}
\label{sec:trajectories}

The IGD trajectories in Fig.~\ref{fig:trajectories} exhibit two recurrent
patterns: \emph{immediate separation} (e.g., WFG4), where IVF-SPEA2 diverges
favorably within the first 10--15\% of generations; and \emph{delayed
	crossover} (e.g., DTLZ3), where SPEA2 briefly leads before IVF-SPEA2
overtakes later. These representative patterns illustrate that operator
benefit can emerge at different phases of the run.

\begin{figure*}[t]
	\centering
	\includegraphics[width=\textwidth]{figures/fig2_trajectories.pdf}
	\caption{IGD trajectories for four representative instances (median and IQR
		over 30 paired seeds): (a)~DTLZ3 delayed crossover; (b)~WFG4 gradual
		divergence; (c)~WFG8 near-neutral; (d)~WFG2 harmful case.
		The vertical dashed line at 20\% of generations marks the
		warmup horizon used by the controller. Trajectories reveal that
		IVF benefit can emerge immediately or after a crossover.}
	\label{fig:trajectories}
\end{figure*}

The collective acceptance criterion (H2) self-limits IVF intensity over time.
Fig.~\ref{fig:cycles} shows the resulting archive turnover for three
representative cases: higher-turnover beneficial cases remain above the
controller threshold ($\theta \approx 0.216$) for longer than neutral
or harmful ones. Early generations typically execute both IVF cycles
($\bar{c} \approx 1.8$), whereas later generations use fewer
($\bar{c} \approx 1.1$).

\begin{figure}[t]
	\centering
	\includegraphics[width=\columnwidth]{figures/fig4_turnover.pdf}
	\caption{Mean archive turnover (IVF-SPEA2) for three representative
		cases. The vertical line marks the 20\% warmup horizon;
		the horizontal dotted line denotes the controller threshold
		$\theta \approx 0.216$. Higher-turnover beneficial cases
		remain above the threshold longer than neutral or harmful ones.}
	\label{fig:cycles}
\end{figure}

\subsection{Early Dynamics as Candidate Discriminants}
\label{sec:early_turnover}

Archive turnover---the fraction of archive members replaced per
generation---reflects how much the population changes between consecutive
generations. The IVF-to-SPEA2 turnover ratio is ${\approx}\,1.0$ across
instances, so the discriminative signal comes from absolute turnover levels
rather than from the ratio. \textsc{Helpful} cases exhibit higher absolute
early turnover than \textsc{Not helpful} cases.

We test 18 case-level summary features with Mann--Whitney $U$
tests~\cite{mann1947test} and BH correction at $q = 0.05$. The set includes
early-run quantities available at decision time (turnover, ND fraction,
fitness, spacing, cycle counts) and endpoint references not available at
decision time (paired IGD/HV deltas).
Table~\ref{tab:dynamic_features} reports the five smallest raw $p$-values.
Both IVF and SPEA2 early turnover survive BH correction, but only the
IVF-side signal is usable without a parallel baseline run.

\begin{table}[t]
	\centering
	\caption{Top five dynamic features by raw $p$-value on the 51 labeled
		cases. Two early-turnover features survive BH correction.}
	\label{tab:dynamic_features}
	\begin{tabular}{lccc}
		\toprule
		\textbf{Dynamic feature}       & $\boldsymbol{p_{\mathrm{raw}}}$ &
		$\boldsymbol{q_{\mathrm{BH}}}$ & $\boldsymbol{\hat{A}_{12}}$                     \\
		\midrule
		Mean turnover (IVF)            & 0.0007                          & 0.013 & 0.849 \\
		Mean turnover (SPEA2)          & 0.0017                          & 0.015 & 0.824 \\
		Final HV delta                 & 0.028                           & 0.169 & 0.727 \\
		Final IGD delta                & 0.073                           & 0.295 & 0.685 \\
		Mean IVF cycles (early)        & 0.098                           & 0.295 & 0.335 \\
		\bottomrule
	\end{tabular}
\end{table}

Among the usable summaries, mean early IVF turnover yields the strongest
signal: $p_{\mathrm{raw}} = 7.2 \times 10^{-4}$,
$q_{\mathrm{BH}} = 0.013$, $\hat{A}_{12} = 0.849$, with class medians
$0.250$ (\textsc{Helpful}) and $0.209$ (\textsc{Not helpful}). A
$10{,}000$-iteration label-permutation test confirms robustness: only 3
permutations reached the observed $\hat{A}_{12}$ (empirical
$p = 3 \times 10^{-4}$; 99th percentile of the null distribution is $0.737$).
Fig.~\ref{fig:turnover_boxplot} shows both the label separation and its
relation to final IGD delta.

A leave-one-out threshold rule on this feature
reaches $\mathrm{BA} = 0.691$ ($\mathrm{MCC} = 0.311$). An oracle rule
using SPEA2 early turnover reaches $0.704$ but requires a parallel baseline
run.

\begin{figure}[t]
	\centering
	\includegraphics[width=\columnwidth]{figures/fig4_discriminant.pdf}
	\caption{Early-turnover signal on the 51 labeled cases.
    \textsc{Helpful} cases exhibit higher early turnover than
    \textsc{Not helpful} cases ($\hat{A}_{12}=0.849$, $q_{\mathrm{BH}}=0.013$),
    with the median decision threshold marked
    ($\theta \approx 0.216$, dotted).}
	\label{fig:turnover_boxplot}
\end{figure}

Under leave-one-family-out refitting, the threshold rule reaches
$\mathrm{BA} = 0.754$ ($\mathrm{MCC} = 0.413$, sensitivity $0.707$,
specificity $0.800$) with median threshold $\approx 0.216$.
The higher BA under leave-one-family-out versus LOOCV
($\mathrm{BA} = 0.691$) arises because thresholds are refitted per held-out
family rather than per held-out instance; given the small and imbalanced
case set (41:10), this estimate should be interpreted as promising but
provisional. The signal is sensitive to the warmup horizon: at 10\% of
generations, IVF turnover remains raw-significant ($p = 0.016$,
$\hat{A}_{12} = 0.749$) but no longer survives BH correction
($q_{\mathrm{BH}} = 0.110$). Fig.~\ref{fig:horizon} summarizes these
trends.

\begin{figure}[t]
	\centering
	\includegraphics[width=\columnwidth]{figures/fig8_warmup_horizon.pdf}
	\caption{Sensitivity of the turnover signal to the warmup horizon.
		Panel~(a): switching performance (balanced accuracy).
		Panel~(b): discriminative strength ($\hat{A}_{12}$).
		The shaded band at 20\% marks the chosen horizon, which
		balances early decision-making ($\hat{A}_{12} = 0.849$,
		$\mathrm{BA} = 0.754$) against stable class separation;
		at 10\% the $\hat{A}_{12}$ signal is weaker (0.749) and
		does not survive BH correction.}
	\label{fig:horizon}
\end{figure}

Early turnover thus separates \textsc{Helpful} from \textsc{Not helpful} cases
more cleanly than static descriptors when the goal is a binary warmup
decision. Static descriptors remain useful for ranking expected gain, but
turnover is the only signal in this study that both survives multiple-testing
correction and supports a simple within-run rule.
However, this signal is family-dependent: ZDT, DTLZ, and MaF
have few \textsc{Not helpful} cases, and in WFG (where both
classes are represented) the separation is weaker
($\hat{A}_{12}=0.775$, $p=0.055$), foreshadowing the controller's
failure mode discussed in Section~\ref{sec:ctrl_results}.


%% ===================================================================
\section{Turnover-Threshold Warmup Controller}
\label{sec:controller}

\begin{figure*}[t]
	\centering
	\includegraphics[width=\textwidth]{figures/fig9_wfg_failure_mode.pdf}
	\caption{Failure mode of the early-turnover controller on WFG configurations.
		Panel~(a) shows the early-turnover score during the first 20\% of the
		evaluation budget, with the deactivation threshold $\theta=0.248$.
		Configurations below the threshold are deactivated by the controller.
		Panel~(b) reports $\Delta\mathrm{HV} = \mathrm{HV}_{\mathrm{CTRL}} -
		\mathrm{HV}_{\mathrm{IVF/SPEA2}}$; negative values indicate a loss
		caused by the controller relative to always-on IVF/SPEA2. Helpful
		configurations below the threshold correspond to premature deactivation
		events and explain the observed controller losses on WFG.}
	\label{fig:wfg_failure}
\end{figure*}

We test a one-shot warmup switch that keeps or drops the IVF operator based
on early archive turnover. Unlike adaptive operator selection schemes, it
makes a single irreversible decision after warmup and does not revisit it.

\subsection{Controller Design}
\label{sec:ctrl_design}

During a warmup phase (first 20\% of generations), the algorithm runs
IVF-SPEA2 and records mean archive turnover $\bar{\tau}$. If
$\bar{\tau} \geq \theta$, IVF remains active; otherwise, the algorithm
reverts to standard SPEA2 for the remaining 80\%. Thresholds are imported
from leave-one-family-out calibration in
Section~\ref{sec:early_turnover} and frozen before execution. For each
held-out family, $\theta$ is estimated from the remaining families by
exhaustive search over candidate thresholds (midpoints between adjacent
unique turnover values) to maximize balanced accuracy; ties are resolved
by maximizing the Matthews correlation coefficient, then accuracy.
This yields $\theta \approx 0.216$ for ZDT, DTLZ, and MaF, and
$\theta \approx 0.248$ for WFG. In deployment, the appropriate threshold
is selected based on known benchmark family membership. For truly unseen
problems where the family is unknown, a single global threshold would be
required; its performance would likely be weaker than the family-wise
results reported here. The controller is a single-feature design with
calibration external to the evaluated family.

When IVF is dropped, the current population and archive are retained without
modification, and only the IVF operator is disabled for the remaining
generations; SPEA2 operators (SBX, polynomial mutation, environmental
selection) continue normally.

The controller runs on the same 51 cases with the same 30 paired seeds as the
dynamics batch, enabling three-way comparison (CTRL vs.\ always-on IVF-SPEA2
vs.\ SPEA2). Statistical comparison uses paired Wilcoxon signed-rank tests
with BH correction applied across the 51 instance-level tests
within each algorithm comparison and endpoint. Case-wise
win/tie/loss counts are reported
conservatively: a case counts as a win or loss only when the corrected paired
test indicates a clear directional difference under the case-level comparison
rule used in the released analysis pipeline; otherwise it is treated as a
tie. Hypervolume is the primary controller endpoint to avoid reusing the
IGD-based labels; IGD is checked as a supporting endpoint.

\subsection{Results}
\label{sec:ctrl_results}

Table~\ref{tab:ctrl_wtl} reports win/tie/loss counts stratified by the
operator-benefit label from Section~\ref{sec:fla_setup} (labels derived from
IGD in the independent-run cohort).

\begin{table}[t]
	\centering
	\caption{Win/tie/loss counts for CTRL against IVF-SPEA2 and SPEA2 under HV
		with family-wise out-of-sample thresholds.}
	\label{tab:ctrl_wtl}
	\begin{tabular}{llccc}
		\toprule
		\textbf{Comparison} & \textbf{Label}           & \textbf{W} & \textbf{T} & \textbf{L} \\
		\midrule
		CTRL vs IVF-SPEA2   & \textsc{Helpful} (41)      & 0          & 32         & 9          \\
		CTRL vs IVF-SPEA2   & \textsc{Not helpful} (10) & 2          & 8          & 0          \\
		\cmidrule(lr){1-5}
		CTRL vs SPEA2       & \textsc{Helpful} (41)      & 33         & 8          & 0          \\
		CTRL vs SPEA2       & \textsc{Not helpful} (10) & 1          & 6          & 3          \\
		\bottomrule
	\end{tabular}
\end{table}

On the 10 \textsc{Not helpful} instances---where always-on IVF degrades or does
not improve SPEA2---the controller never loses to IVF-SPEA2 (2 wins, 8 ties,
0 losses). On the 41 \textsc{Helpful} instances, it ties IVF-SPEA2 in 32 cases
but loses in 9, all in WFG. The main failure mode is premature
deactivation on that family. On WFG instances, both \textsc{Helpful}
(mean turnover $0.218$) and \textsc{Not helpful} (mean $0.196$) cases fall
below the WFG-specific threshold ($\theta \approx 0.248$), so the
controller deactivates IVF for all 18~WFG instances. Although
\textsc{Helpful} cases tend toward higher turnover
($\hat{A}_{12}=0.775$, $p=0.055$), the class separation is insufficient
to avoid premature deactivation under a fixed threshold. A plausible
explanation is that WFG's multi-level deceptive structure sustains high
archive churn even when convergence stalls, which weakens turnover as a
proxy for search progress (Fig.~\ref{fig:wfg_failure}).
This interpretation is consistent with the concentration of controller
losses in that family, but it remains a mechanistic hypothesis;
establishing causality would require controlled trajectory-level
intervention experiments.

Relative to SPEA2, the controller records 34 wins, 14 ties, and 3 losses
in hypervolume (33 of the wins on \textsc{Helpful} cases). On the
10~\textsc{Not helpful} cases, the controller does not uniformly protect
against SPEA2 (1~win, 6~ties, 3~losses). The IGD summary follows the same
direction but is less pronounced: 41/51 match-or-beat versus IVF-SPEA2 and
47/51 versus SPEA2. Because the label distribution is imbalanced toward
always-on IVF (41/51 \textsc{Helpful}), raw label frequency alone is not the
right target. We thus read the controller as a risk-management mechanism:
it avoids hypervolume losses against always-on IVF on all 10 instances where
IVF is unhelpful, at the cost of forfeiting some gains on 9 of 41 beneficial
instances---all concentrated in WFG. The controller should therefore be
interpreted as a risk-management mechanism relative to always-on IVF, not as a
guarantee of recovery to canonical SPEA2 performance.
Table~\ref{tab:ctrl_wtl} summarizes the controller performance under
family-wise out-of-sample thresholds; Fig.~\ref{fig:wfg_failure} details
the concentration of losses in that family.

%% ===================================================================
\section{Discussion and Conclusions}
\label{sec:conclusion}

Taken together, the evidence supports a focused claim: for one-shot operator
switching inside SPEA2, warmup dynamics are more actionable than ex ante
landscape descriptors. Static ELA-style features help rank expected gain
magnitude, but this signal does not survive cross-family binary switching.
Early turnover over the first 20\% of generations separates
\textsc{Helpful} from \textsc{Not helpful} cases more reliably and supports a
one-shot warmup rule. Under family-wise out-of-sample calibration, the rule
matches or beats always-on IVF-SPEA2 on 42/51 hypervolume cases and incurs
zero hypervolume losses on the 10~\textsc{Not helpful} cases. The trade-off is
9~losses on \textsc{Helpful} cases, all in WFG, where turnover is a weak proxy
for convergence progress. The controller is better interpreted as
risk reduction in a fixed host than as universal protection.

These findings are specific to one operator in one host algorithm, but they
address a practical design decision: whether to keep an expensive
intensification module after warmup. In this setting, static descriptors are
more useful as priors on expected gain magnitude, while early dynamics are
more suitable for the actual keep/drop decision.

\subsection{Limitations}
Our FLA analysis uses ELA-style features; more specialized descriptors such as
Pareto local optima networks~\cite{ochoa2024funnels} or learned features
(Deep-ELA~\cite{seiler2024deep}) remain untested. The sample size
($n = 51$) is modest and class-imbalanced (41:10), which makes threshold
calibration uncertain. The controller uses a single feature with family-wise
thresholds imported from leave-one-family-out analysis---stronger than
in-sample calibration, but not full unseen-domain validation. All
9~hypervolume losses to IVF-SPEA2 on \textsc{Helpful} cases occur in WFG
(Section~\ref{sec:ctrl_results}); the near-overlap of turnover distributions
on that family limits discriminative power. The 9~losses on \textsc{Helpful}
and 3~losses to SPEA2 on \textsc{Not helpful} leave engineering margin.
Controller results are endpoint-specific: we use hypervolume to avoid reusing
IGD-based labels, and the IGD summaries are slightly less favorable. The
study examines one operator (IVF) in one host (SPEA2); broader
generalization remains open. In addition, raw-objective landscape
descriptors may partially encode benchmark-scale differences, especially
for geometry features such as front diameter.

\subsection{Future work}
Future work should first test whether the turnover signal transfers to
unseen benchmark suites and real engineering problems. A second direction
is to replace the single-threshold rule with multi-feature controllers
that combine turnover with progress proxies, such as early HV gain,
archive stability, or dominance-depth changes. Controlled
trajectory-level interventions are needed to test whether the WFG failure
mode arises from deceptive archive churn rather than useful convergence
progress. Richer static descriptors (e.g., PLON or Deep-ELA) may also
strengthen the prior-signal role of landscape features.

\subsection{Data and code availability}
Code, processed data, manuscript sources, and reproducibility artifacts for
this study are available at
\url{https://github.com/pedsanches/IVF-SPEA2/releases/tag/submission-snapshot-2026-03}.
Additional figures (family-wise turnover separability, IGD-HV consistency,
and experimental design schematic) are included in the supplementary
material.

\bibliographystyle{IEEEtran}
\bibliography{references}

\end{document}
